\section{기호와 manifest 식별자}
\label{app:symbols}
\begin{longtable}{@{}P{.20\linewidth}P{.74\linewidth}@{}}
\caption{주요 기호의 의미}\\
\toprule 기호 & 의미 \\\midrule
\endfirsthead
\toprule 기호 & 의미 \\\midrule
\endhead
\bottomrule\endfoot
$D$ & 표본의 수치 표현 \\
$T_{\mathrm{val}}$ & 벡터ㆍepochㆍ행렬ㆍ사건을 포함한 값 형식 \\
$\mathcal C,\chi_c$ & 순서가 있는 채널 스키마와 채널 기술자 \\
$\kappa,\mathcal R$ & 시계 영역과 표본 발생 규칙 \\
$X\in D^{C\times N}$ & 채널 $C$개, 표본 $N$개의 값 블록 \\
$I_i^{\mathrm{meas}}$ & 대표 측정 시각의 구간 \\
$a_i$ & 지정한 관측 경계에서의 가용 시각 \\
$\cut_u(S)$ & 시각 $u$까지 가용한 스트림 부분 \\
$\rho_i,\rho_v^{\mathrm{fit}}$ & 레코드 계보와 모델 학습 계보 \\
$e^{\mathrm{schema}},e^{\mathrm{clock}}$ & 스키마와 시계의 세대 식별자 \\
$H,\theta,\Omega_H$ & 하드웨어 프로필, 선택 구성, 유효 구성 공간 \\
$C(H,\theta),R_P$ & 구성별 기능 집합과 프로그램의 기능 요구 \\
$Z_o,\mathcal J_o,\delta_o$ & 연산자 상태 공간, 상태 불변식, 상태 전이 \\
$\mathcal E_o$ & 처리 모드, 요구 기능 집합, 전방 관측 길이 \\
$M_v,\vartheta_v$ & 모델 버전과 수치 매개변수 \\
$\alpha,\beta$ & 누적 유입과 서비스의 경계 함수 \\
$b,r,R,T$ & burst, 지속 유입률, 서비스율, 서비스 지연 \\
\end{longtable}

manifest의 모든 식별자는 이름 공간과 버전을 함께 가진다.
장치ㆍ수신기 식별자, 원천 스트림 식별자, 획득 epoch, 데이터 블록 순번은
원천 레코드의 범위를 결정한다. 스키마ㆍ시간 뷰ㆍ모델ㆍ계약 식별자는
변환 관계를 결정한다. 수치 상수와 설정 레코드는 정규화한 표현의 digest를 사용한다.
기록기는 이 식별자와 전이 사건을 함께 저장하여 동일 원천에서 생성한 여러 분석 결과를 연결한다.
